% ATTEMPT_STATUS: unresolved
% RESEARCH_GENERATED: 2026-10-01; GPT-6 Astra Ultra; individual mathematical attempt
% TOP_PROBLEM: {"id": 150, "title": "Bounds for Birch's Theorem", "category_id": 1, "category": {"id": 1, "name": "number_theory", "display_name": "Number Theory", "description": "Properties of integers, prime numbers, Diophantine equations.", "slug": "number-theory", "order_index": 1, "created_at": "2026-07-31T15:26:25.670Z"}, "status": "open"}
% FIRST_POSTED: 2026-10-01
% Imported from: unsolvedmath_additional_500.tex; UnsolvedMath ID 150
% !TeX program = lualatex
% Compile twice: lualatex 150.tex
% Self-contained: no external bibliography, images, or \input files.
% Numeric section labels and visible numbers are original UnsolvedMath IDs.
\documentclass[11pt,a4paper]{article}
\usepackage[margin=24mm,headheight=14pt]{geometry}
\usepackage{fontspec}
\setmainfont{Latin Modern Roman}
\setsansfont{Latin Modern Sans}
\setmonofont{Latin Modern Mono}
\usepackage{amsmath,amssymb,mathtools}
\usepackage{unicode-math}
\setmathfont{Latin Modern Math}
\usepackage{microtype}
\usepackage{enumitem,xurl,needspace}
\usepackage[dvipsnames]{xcolor}
\usepackage{hyperref}
\hypersetup{unicode=true,colorlinks=true,linkcolor=MidnightBlue,urlcolor=MidnightBlue,
 pdftitle={UnsolvedMath Problem 150},pdfauthor={Source-annotated compilation},bookmarksnumbered=true}
\usepackage{bookmark,tocloft,titlesec,fancyhdr}
\setlength{\cftsecnumwidth}{4.2em}
\cftsetpnumwidth{2.5em}
\cftsetrmarg{4em}
\titleformat{\section}{\Large\bfseries\raggedright}{\thesection}{0.7em}{}
\pagestyle{fancy}\fancyhf{}
\fancyhead[L]{\small\sffamily UnsolvedMath research attempt}
\fancyhead[R]{\small Original catalogue IDs}
\fancyfoot[C]{\thepage}
\setlength{\parindent}{0pt}
\setlength{\parskip}{5pt plus 1pt}
\setlength{\emergencystretch}{3em}
\setcounter{tocdepth}{1}\setcounter{secnumdepth}{1}
\setlist[itemize]{leftmargin=3em,itemsep=4pt,topsep=4pt,parsep=0pt}
\allowdisplaybreaks
\newcommand{\N}{\mathbb{N}}
\newcommand{\Z}{\mathbb{Z}}
\newcommand{\Q}{\mathbb{Q}}
\newcommand{\R}{\mathbb{R}}
\newcommand{\C}{\mathbb{C}}
\newcommand{\F}{\mathbb{F}}
\newcommand{\A}{\mathbb{A}}
\newcommand{\PP}{\mathbb{P}}
\newcommand{\E}{\mathbb{E}}
\newcommand{\eps}{\varepsilon}
\newcommand{\dd}{\,\mathrm d}
\newcommand{\sref}[2]{\hyperlink{source:#1:#2}{[S#2]}}
\newif\ifshowcatalogue
\showcataloguetrue
\newcommand{\cataloguescope}[1]{\ifshowcatalogue
 \paragraph{Statement recorded in the supplied source.}
 #1\fi}
\begin{document}
% UnsolvedMath ID 150; source catalogue code GREEN-062

\setcounter{section}{149}

\section{Optimal Dimension Bounds for Odd-Degree Forms}\label{150}

{\small\sffamily UnsolvedMath ID 150 · Number Theory}

\subsection{Definitions and mathematical statement}

\paragraph{Shared notation.}Unless a section says otherwise, $\N=\{1,2,\ldots\}$,
$\Z$ is the ring of integers, $\Q,\R,\C$ are the rational, real, and complex
fields, $[N]=\{1,\ldots,\lfloor N\rfloor\}$, and $\log$ is the natural
logarithm. For real functions of a parameter tending to infinity,
$f\ll g$ and $f=O(g)$ mean $|f|\leq Cg$ eventually for some fixed $C$;
$f\gg g$ reverses this comparison; $f=o(g)$ means $f/g\to0$;
$f\sim g$ means $f/g\to1$; and $f\asymp g$ means both order bounds.
A subscript on $O$ or $\ll$ specifies allowed constant dependence.
The expression $n^{o(1)}$ in an upper bound means growth smaller than
$n^\epsilon$ for each fixed $\epsilon>0$. Local definitions govern any
exception, finite-field convention, or use of the same symbol for another
object. An asymptotic claim is not automatically an effective algorithm.



A degree-$d$ form is a homogeneous polynomial $F$ satisfying $F(t\mathbf x)=t^dF(\mathbf x)$. A nontrivial integer zero is $\mathbf x\in\Z^n\setminus\{0\}$ with $F(\mathbf x)=0$. For odd $d\geq3$, the source asks for effective or sharp bounds on a dimension threshold valid uniformly over all integer coefficients. Determining the least possible threshold is the corresponding extremal formulation. \sref{150}{1} \sref{150}{2}

\paragraph{Statement recorded in the supplied source.}
Let $d \geq 3$ be odd. Give bounds on $\nu(d)$ such that if $n > \nu(d)$ then any homogeneous polynomial $F(\mathbf{x}) \in \mathbb{Z}[x_1, \dots, x_n]$ of degree $d$ has a nontrivial integer zero. \sref{150}{1}

\paragraph{Residual task reported by the archive.}

The embedded review (2026-08-17) records: Record the best explicit upper and lower bounds for the exact single-form threshold nu(d). \sref{150}{1}

\subsection{Short English statement}

How many variables are needed to guarantee a nonzero integer solution to every homogeneous equation of a fixed odd degree?

\subsection{Research attempt}

\paragraph{Provenance and scope.}
This individual attempt was generated by GPT-6 Astra Ultra on 1 October 2026. The method was to derive explicit partial results and test the first proposed extension adversarially. The original statement and sources below are retained. No claim of mathematical novelty or complete solution is made.

\paragraph{Formal target and updated source scope.}
For odd $d\geq3$, the threshold must work for every integer homogeneous degree-$d$ form, with no restriction on coefficients. A nonzero rational zero can be cleared of denominators, so rational and integer solvability are equivalent here. The inspected Lampert--Snowden--Ziegler manuscript gives polynomial dependence on the number of forms at fixed degree; that is not a degree-uniform estimate for the single-form threshold as $d$ varies. We construct an explicit local obstruction showing that at least $d^2$ variables may be required before a universal theorem can begin.

\paragraph{A norm form with controlled valuations.}
Fix a prime $p$, and choose a monic degree-$d$ polynomial $g(T)\in\Z[T]$ whose reduction modulo $p$ is irreducible. Such polynomials exist because the finite field of order $p^d$ has a degree-$d$ element over $\F_p$. Let $R=\Z[T]/(g)$, free of rank $d$ with basis $1,T,\ldots,T^{d-1}$. For $u=\sum_{i=0}^{d-1}u_iT^i$, define $N(u)$ to be the determinant of multiplication by $u$ on that basis. It is a homogeneous degree-$d$ polynomial with integer coefficients.

If not all coordinates $u_i$ are divisible by $p$, then the reduction of $u$ is a nonzero element of the field $R/pR\cong\F_{p^d}$. Multiplication by it is invertible, so its determinant is nonzero modulo $p$. Thus $p\nmid N(u)$. For a nonzero integer vector $u$, let $v=\min_i v_p(u_i)$. Homogeneity and the preceding observation give
\[
 v_p(N(u))=dv.
\]
This establishes the valuation formula without any appeal to an unproved statement about arbitrary norm forms or ramified extensions.

\paragraph{A degree-$d$ form in $d^2$ variables with no nonzero zero.}
Take $d$ independent blocks $u^{(0)},\ldots,u^{(d-1)}$, each of length $d$, and define
\[
 F=\sum_{j=0}^{d-1}p^j N(u^{(j)}).
\]
It is an integer homogeneous form of degree $d$ in $d^2$ variables. If a block is nonzero, its corresponding summand has valuation congruent to $j$ modulo $d$. Therefore the nonzero summands have pairwise distinct valuations. In any nonempty finite collection of integers with a unique lowest $p$-adic valuation, their sum cannot vanish: divide by the lowest power of $p$, and reduction modulo $p$ leaves exactly one nonzero term. Hence $F=0$ forces every block to be zero.

We have proved that the least universal threshold satisfies $\nu(d)\geq d^2$ under the convention that every form in more than $\nu(d)$ variables has a nontrivial zero. The construction actually works for every positive degree; oddness is not needed for this local obstruction. For odd degree it shows why the automatic real sign change is far from sufficient for rational solvability.

\paragraph{A concrete cubic instance.}
For $d=3$ and $p=2$, use $g(T)=T^3-T-1$, whose reduction $T^3+T+1$ has no root in $\F_2$ and is therefore irreducible. Multiplication by $a+bT+cT^2$ has matrix
\[
 M(a,b,c)=
 \begin{pmatrix}
 a&c&b\\ b&a+c&b+c\\ c&b&a+c
 \end{pmatrix}.
\]
Each column follows by multiplying the corresponding basis element and reducing $T^3=T+1$. Thus the explicit nine-variable cubic
\[
 \det M(u^{(0)})+2\det M(u^{(1)})+4\det M(u^{(2)})
\]
has no nonzero integer zero. The determinant expression itself is a fully specified polynomial and avoids transcription errors in an expanded cubic. The valuation residues are zero, one, and two modulo three, so the noncancellation proof applies verbatim.

\paragraph{Why a real-zero argument cannot bypass the obstruction.}
For odd degree, $F(-x)=-F(x)$ and continuity on a sphere can produce real zeros when there are enough directions. Approximating one by rational points makes $F$ small, not zero. After clearing denominators, the small error is multiplied by a degree-$d$ power and need not disappear. The explicit form above supplies an adversarial example: it has no nonzero rational zero despite the real sign symmetry. Any argument using only continuity and rational density is therefore missing an arithmetic step.

\paragraph{Verification and remaining gap.}
The norm determinant is integral, homogeneous, and its reduction is invertible exactly when the residue-field element is nonzero. Zero blocks are omitted before choosing the minimum valuation, avoiding an undefined valuation of zero. The cubic reduction has been checked at both field elements, and the matrix columns are displayed. The result is a rigorous quadratic lower obstruction, not an upper bound or an optimality assertion. The unresolved task is to obtain effective upper bounds with carefully tracked dependence on degree, and ultimately to determine how close the true threshold can be to this $d^2$ barrier.

\subsection{Sources}

{\small

\begin{itemize}

\item[\hypertarget{source:150:1}{S1}] ULAM.AI, \emph{UnsolvedMath}, supplied \texttt{problems.json}, record 150 (GREEN-062), statement and background. The embedded literature-review date is 2026-08-17; that is the archive’s review date, not a fresh certification. Dataset landing page: \url{https://huggingface.co/datasets/ulamai/UnsolvedMath}.

\item[\hypertarget{source:150:2}{S2}] Ben Green, 100 Open Problems (author’s maintained problem list). \url{https://people.maths.ox.ac.uk/greenbj/papers/open-problems.pdf}. The author-hosted document was inspected on 27 September 2026; the archive’s local code is not a problem-number locator in this paper.

\item[\hypertarget{source:150:3}{S3}] B. J. Birch, Homogeneous forms of odd degree in a large number of variables, Mathematika 4 (1957). \url{https://doi.org/10.1112/S0025579300001145}. Bibliographic information imported from the supplied record.

\item[\hypertarget{source:150:4}{S4}] A. Lampert, A. Snowden, T. Ziegler, Polynomial Bounds for Birch's Theorem, arXiv:2512.00697 (2025). \url{https://arxiv.org/abs/2512.00697}. Bibliographic information imported from the supplied record.

\item[E1] A. Lampert, A. Snowden and T. Ziegler, \emph{Polynomial Bounds for Birch's Theorem}, 2025. \url{https://arxiv.org/abs/2512.00697}. Inspected 1 October 2026.

\end{itemize}

}

\label{end:150}
\end{document}
